\section{研究目标与实验思路}

本研究围绕中文对联生成，完成字符级GPT搭建、古典语料预训练、任务微调与推理展示。实现路线参考Raschka的\emph{Build a Large Language Model (From Scratch)}及其\href{https://github.com/rasbt/LLMs-from-scratch}{配套资源}\cite{raschka}，将词元表示、因果注意力和下一词元预测落实为可训练的小参数模型。

原始GPT接受对联配对监督后，已能生成五字或七字下联，但仍出现同位照用、无依据重复和词组对应较弱等问题。为此，本文在自建主干上引入ChineseBERT条件表示、词汇注意力与位置对应连接，并比较无规则生成和规则约束生成。另从同一预训练底座独立微调古诗文赏析模型，观察监督学习对输出任务和体裁的改变。

全文先说明语言建模机制及基础模型实现，再介绍预训练数据、训练过程与迁移结果；对联部分按照方法、结果和讨论组织，最后总结典型局限与后续研究方向。
